\section{Case Studies}\label{sec:usecases}

In this section we conclude our case studies. See \cref{sec:eval} for the setup
and methodology behind each study.


% ===========================================================
\subsection{``Battery Saver'' Mode}\label{subsec:usecase-battery}

% About: background
In this first case study, we revisit the motivational Forward switching example
(\cref{sec:motivation}). Our goal is to complete a workload within a strict
energy budget. Here, our netbook's storage will use {\em Freestyle Balanced}
when operating normally and {\em ChaCha8} when battery saver mode is activated.

% About: setup
To simulate this, we repeat the following twice: 1) we begin writing 10 40MB
files using Freestyle Balanced. 2) After five seconds, our device simulates
battery saver mode by pinning \sys processes to the energy-efficient LITTLE
cores and underclocking them to lowest frequency. 3) We complete the workload.

\def \hmina {\hspace{-0.1in}}
\def \hminb {\hspace{-0.2in}}

\def \fgw {2in}
\def \fgh {1in}

% \begin{floatingfigure}[r]{2in}  (and \end{..})

\begin{figure}[t]
    \centerline{
        % \hmina
        % \includegraphics[width=\fgw]{data/usecase-battery}
        % \hminb
        \includegraphics[height=\fgh]{empty.eps}
    }

    \mycaption{fig:usecase-battery}{Battery saver case}{Energy-security
    tradeoff given strict energy budget as discussed in
    \cref{subsec:usecase-battery}.}
\end{figure}



On the first run, the workload is completed normally. On the second run, when
entering battery saver mode, we additionally trigger a switch to the ChaCha8
crypt using the Forward switching model.

% About: outcome
\Cref{fig:usecase-battery} shows the total energy used after completing both
runs. {\em FB only} (left) favors security even when backups are paused while
{\em Forward (FB+C8)} (right) performs the switch when the system enters battery
saver mode. These results show that using Forward switching resulted in a 3.3x
total energy reduction compared to exclusively using Freestyle Balanced,
allowing us to remain within our energy budget.


% ===============================
\subsection{No-Downtime Encryption Upgrade}\label{subsec:usecase-upgrade}

% About: background
In this second case study we revisit the Mirrored switching example
(\cref{subsec:des-switch}). Our goal is to upgrade live storage from one cipher
to another without downtime or degraded service. Our operator wants to upgrade
from {\em ChaCha20} to {\em Freestyle Fast}.

% About: setup
To simulate pre- and post-migration state, we execute 10 5MB write-read
operations to two \sys instances using ChaCha20 and Freestyle Fast respectively.
To simulate the activity during migration, we repeat the operations on a third
instance using Mirrored switching.

\def \hmina {\hspace{-0.1in}}
\def \hminb {\hspace{-0.2in}}

\def \fgw {2in}
\def \fgh {1in}

% \begin{floatingfigure}[r]{2in}  (and \end{..})

\begin{figure}[t]
    \centerline{
        % \hmina
        % \includegraphics[width=\fgw]{usecase-upgrade}
        % \hminb
        \includegraphics[height=\fgh]{empty.eps}
    }

    \mycaption{fig:usecase-upgrade}{Cipher upgrade without downtime}{Switching
    storage system encryption with zero downtime as discussed in
    \cref{subsec:usecase-upgrade}.}
\end{figure}


% About: outcome
\Cref{fig:usecase-upgrade} shows the overhead of Mirrored switching on our
system's performance. On the left we see similar read latency between pre-,
during-, and post-migration workloads of \TODO{X}s. On the right, we see write
latency increase by \TODO{X}s during the migration but recover post-migration.
Thus, Mirrored switching allows us to transition our storage from one cipher to
another without interruption or egregious overhead.


% ================================================================
\subsection{Scalable Encryption}\label{subsec:usecase-agnostic}

% About: background
In this final case study we revisit the Selective switching example
(\cref{subsec:des-switch}). Our goal is to achieve a performance win by ensuring
that the only data encrypted using the strongest highest-overhead crypt is the
small percentage of data that needs it. We note the Selective switching model
requires users to annotate certain files with a special tag via file system
calls and stored in the inode. Because the \sys layer is file-oblivious, every
block I/O through the \sys layer will be labeled with the corresponding crypt.

% About: setup
We begin with 10 5MB and 10 4KB write-read operations to two \sys instances: one
using {\em ChaCha8}, a high performance crypt, and the other using {\em
Freestyle Strong}, a high security crypt. We repeat the operations on a third
instance using Selective switching where 30\% of the data is considered highly
sensitive and uses Freestyle Strong.

\def \hmina {\hspace{-0.1in}}
\def \hminb {\hspace{-0.2in}}

\def \fgw {2in}
\def \fgh {1in}

% \begin{floatingfigure}[r]{2in}  (and \end{..})

\begin{figure}[t]
    \centerline{
        % \hmina
        % \includegraphics[width=\fgw]{data/usecase-agnostic}
        % \hminb
        \includegraphics[height=\fgh]{empty.eps}
    }

    \mycaption{fig:usecase-agnostic}{Scalable encryption}{Filesystem-agnostic
    file-level encryption as discussed in \cref{subsec:usecase-agnostic}.}
\end{figure}


% About: outcome
{\em C8 only} (leftmost) and {\em FS only} (rightmost) in
\cref{fig:usecase-agnostic} show the two extremes: ChaCha8 exhibits low latency,
consistently less than 1s across workloads, while Freestyle Strong exhibits 5s
to 20s latency depending on the workload. {\em Selective (C8+FS)} (middle) shows
a reduction of 3.1x to 4.8x for read latency and 1.6x to 2.8x for write latency,
all without compromising on security. Thus, Selective switching keeps our
sensitive data at the mandated security level while keeping the performance
benefits of using a fast crypt.
